\documentclass[10pt,a4paper]{amsart}
\usepackage[margin=19mm]{geometry}
\usepackage{amsmath,amssymb,mathtools}
\usepackage[scr=rsfs,cal=euler]{mathalfa}
\usepackage{xfrac}
\usepackage[unicode,hidelinks,bookmarks=false]{hyperref}
\newcommand{\ie}{i.e.\ }
\newcommand{\cf}{cf.\ }
\newcommand{\resp}{respectively\ }
\newcommand{\Subsection}{\subsection}
\providecommand{\nobreakdash}{\nobreak\hbox{-}}
\setlength{\emergencystretch}{2em}
\multlinegap0pt
\usepackage{bm}
\usepackage{bbm}
\usepackage{dsfont}
\usepackage{xspace}
\usepackage{cancel}
\usepackage{mathtools}
\usepackage[shortlabels]{enumitem}
\usepackage{amsthm}
\makeatletter
\@ifundefined{theorem}{\newtheorem{theorem}{Theorem}}{}
\@ifundefined{definition}{\newtheorem{definition}[theorem]{Definition}}{}
\@ifundefined{lemma}{\newtheorem{lemma}[theorem]{Lemma}}{}
\makeatother
\def\a{\alpha}
\def\b{\beta}
\def\g{\gamma}
\def\v{\varepsilon}
\title[Post-Hopf group algebras, Hopf group braces and Rota-Baxter ]{Post-Hopf group algebras, Hopf group braces and Rota-Baxter operators on Hopf group algebras}
\author{Yan Ning, Xing Wang, Daowei Lu}
\date{Source: arXiv:2507.20176v2 (CC BY 4.0)}
\begin{document}
\maketitle

\noindent\emph{Source and licence.} Post-Hopf group algebras, Hopf group braces and Rota-Baxter operators on Hopf group algebras. Version arXiv:2507.20176v2 (CC BY 4.0), distributed under the Creative Commons Attribution 4.0 International licence, as recorded in the arXiv licence metadata for this exact version and shown on the abstract page https://arxiv.org/abs/2507.20176v2. Full rights notice: licenses/arxiv-2507.20176-CC-BY-4.0.txt.\par\smallskip

\noindent\textit{Editorial scope.} Counted unit: the Lemma whose source label is \texttt{lem:nH} in the article (source0.tex lines 860--878, source label \texttt{lem:nH}), delivered together with the article's own complete original proof of it (source0.tex lines 879--946, one \texttt{proof} environment of 68 lines). No mathematics is written, repaired or re-derived: statement and proof are the article's own text line for line. Required context retained uncounted: eq:RBO (lines 788--794). Every definition, display and label of those lines is retained unchanged. Omitted from this package and not redistributed: 1--787, 795--859, 947--986. Nothing inside a retained excerpt is omitted or altered: every retained body is delivered exactly as the article's lines are written, so no omission receipt of any kind accompanies this package. Citations: the retained excerpts carry no citation command at all, so no bibliography entry of the article is transcribed and no reference is redistributed. Reference adaptation: none was needed and none was made; every reference inside the delivered unit resolves inside it, so no unresolved reference or citation is produced. Printed slips kept literal: no printed word, symbol, hypothesis or number of the counted statement or of its proof was altered. Layout and class: the article's own class is \texttt{article} with packages including amscd, amsfonts, amsmath, amssymb, amsthm, bbm, enumerate, fancyhdr, graphicx, hyperref, indentfirst, latexsym, mathrsfs, babel; the wrapper is the lane's pinned \texttt{amsart} editorial wrapper at 10pt on A4, which restates the theorem-like environments the retained text needs and adds the article's own macro definitions that the retained text needs (\texttt{\textbackslash a}, \texttt{\textbackslash b}, \texttt{\textbackslash g}, \texttt{\textbackslash v}), with extra wrapper packages (bm, bbm, dsfont, xspace, cancel, mathtools, enumitem). The delivered numbering is the unit's own. Assets: no retained span contains a figure environment, a graphics-inclusion command, a PDF-page inclusion, a table or a TikZ picture; no image file is copied into this package, the package declares no asset dependency and no image file is redistributed with it. Licence: version arXiv:2507.20176v2 of this article is distributed under the Creative Commons Attribution 4.0 International licence, as recorded in the arXiv licence metadata for this exact version and shown on the abstract page https://arxiv.org/abs/2507.20176v2. A full rights notice is delivered at licenses/arxiv-2507.20176-CC-BY-4.0.txt. Characters: every retained excerpt is pure ASCII, so the delivered engine, XeLaTeX with its default Latin Modern text font, reports no missing character.
\begin{definition}
Let $H=\{H_\a\}_{\a\in\pi}$ be a cocommutative Hopf $\pi$-algebra. The family of coalgebra homomorphisms $B=\{B_\a:H_\a\rightarrow H_{\bar{\a}}\}_{\a\in\pi}$ is called a Rota-Baxter operator on $H$ if for all $\a,\b\in\pi, a\in H_\a,b\in H_\b$,
\begin{equation} \label{eq:RBO}
  B_\a(a)B_\b(b)=B_{\b\a}\big(a_{(1,\a)}B_\a(a_{(2,\a)})bS_{\bar{\a}}(B_\a(a_{(3,\a)}))\big).
\end{equation}
The pair $(H,B)$ is called a Rota-Baxter Hopf $\pi$-algebra.
\end{definition}

\begin{lemma}\label{lem:nH}
Let $(H=\{H_\a\}_{\a\in\pi},B)$ be a Rota-Baxter Hopf $\pi$-algebra. Define
\begin{eqnarray}
& m_{\a,\b}(g,h)=g\circ_B h:=g_{(1,\a)}B_\a(g_{(2,\a)})hS_{\bar{\a}}(B_\a(g_{(3,\a)})), & \label{eq:RBHm}\\
& T_\a(g):=S_{\bar{\a}}(B_\a(g_{(1,\a)}))S_\a(g_{(2,\a)})B_\a(g_{(3,\a)}), & \label{eq:RBHS}
\end{eqnarray}
for all $g\in H_\a,h\in H_\b, ~\a,\b\in\pi$.
\begin{enumerate}[\quad\rm(1)]
\item $H_B:=(H,m=\{m_{\a,\b}\}_{\a\in\pi},T=\{T_{\a,\b}\}_{\a\in\pi})$ is a cocommutative Hopf $\pi$-algebra. %, called the {\bf descendent Hopf $\pi$-algebra} of $H$.
\item For all $h\in H_\a$, $\a\in\pi$, we have
\begin{eqnarray}
 & B_\a(h_{(1,\a)})B_{\bar{\a}}(T_\a(h_{(2,\a)}))=\v_\a(h)1_e, \label{eq:t} & \\
 & B_{\bar{\a}}T_\a = S_{\bar{\a}}B_\a. &
\end{eqnarray}

\item The operator $B$ is also a Rota-Baxter operator on $H_B$, that is, $(H_B,B)$ is a Rota-Baxter Hopf $\pi$-algebra.
\item The family of maps $B:H_B\rightarrow H$ is a homomorphism of Rota-Baxter Hopf $\pi$-algebras.
\end{enumerate}
\end{lemma}

\begin{proof}
(1) First, we prove $m=\{m_{\a,\b}\}_{\a,\b\in\pi}$ is  associative. For all $g\in H_\a, h\in H_\b, \ell\in H_\g, ~\a,\b,\g\in\pi$,
\begin{eqnarray*}
% \nonumber to remove numbering (before each equation)
& & (g\circ_B h) \circ_B\ell \\
&=& (g_{(1,\a)}B_\a(g_{(2,\a)})hS_{\bar{\a}}(B_\a(g_{(3,\a)}))) \circ_B\ell \\
&=& g_{(1,\a)}B_\a(g_{(2,\a)})h_{(1,\b)}S_{\bar{\a}}(B_\a(g_{(3,\a)}))  B_{\b\a}\big( g_{(4,\a)} B_\a(g_{(5,\a)})h_{(2,\b)}S_{\bar{\a}}(B_\a(g_{(6,\a)})) \big) \\
&&\quad \ell S_{\bar{\b\a}}\big(B_{\b\a}( g_{(7,\a)} B_\a(g_{(8,\a)})h_{(3,\b)}S_{\bar{\a}}(B_\a(g_{(9,\a)})) )\big) \\
&=& g_{(1,\a)}B_\a(g_{(2,\a)})h_{(1,\b)}S_{\bar{\a}}(B_\a(g_{(3,\a)}))  B_\a(g_{(4,\a)})   B_\b(h_{(2,\b)})\ell  S_{\bar{\b\a}}( B_\a(g_{(3,\a)}) B_\b(h_{(3,\b)}) ) \\
&=& g_{(1,\a)}B_\a(g_{(2,\a)})h_{(1,\b)} B_\b(h_{(2,\b)})\ell S_{\bar{\b}}(B_\b(h_{(3,\b)})) S_{\bar{\a}}(B_\a(g_{(3,\a)})) \\
&=& g\circ_B( h_{(1,\b)} B_\b(h_{(2,\b)})\ell S_{\bar{\b}}(B_\b(h_{(3,\b)})) ) \\
&=& g\circ_B(h\circ_B \ell).
\end{eqnarray*}



  Then, we prove $m=\{m_{\a,\b}\}_{\a,\b\in\pi}$ is a $\pi$-coalgebra homomorphism. For all $g\in H_\a, h\in H_\b, ~\a,\b\in\pi$,
\begin{eqnarray*}
& & (g\circ_B h)_{(1,\a\b)} \otimes (g\circ_B h)_{(2,\a\b)} = (g\circ_B h)_{(1,\b\a)} \otimes (g\circ_B h)_{(2,\b\a)} \\
&=& (g_{(1,\a)}B_\a(g_{(2,\a)})hS_{\bar{\a}}(B_\a(g_{(3,\a)})))_{(1,\b\a)} \otimes (g_{(1,\a)}B_\a(g_{(2,\a)})hS_{\bar{\a}}(B_\a(g_{(3,\a)})))_{(2,\b\a)}  \\
&=& g_{(1,\a)}B_\a(g_{(2,\a)})h_{(1,\b)}S_{\bar{\a}}(B_\a(g_{(3,\a)})) \otimes g_{(4,\a)}B_\a(g_{(5,\a)})h_{(2,\b)}S_{\bar{\a}}(B_\a(g_{(6,\a)}))  \\
&=& (g_{(1,\a)}\circ_B h_{(1,\b)})\otimes(g_{(2,\a)}\circ_B h_{(2,\b)}).
\end{eqnarray*}

  Last, we prove $T=\{T_\a\}_{\a\in\pi}$ is an antipode. For all $h\in H_\a, ~\a\in\pi$,
\begin{eqnarray*}
& & m_{\a,\bar{\a}}(id_{H_{\a}} \otimes T_\a)\Delta_\a(h) \\
&=& h_{(1,\a)} \circ_B S_{\bar{\a}}(B_\a(h_{(2,\a)}))S_\a(h_{(3,\a)})B_\a(h_{(4,\a)}) \\
&=& h_{(1,\a)}B_\a(h_{(2,\a)}) S_{\bar{\a}}(B_\a(h_{(3,\a)}))S_\a(h_{(4,\a)})B_\a(h_{(5,\a)}) S_{\bar{\a}}(B_\a(h_{(6,\a)})) \\
&=& h_{(1,\a)} \v_\a(h_{(2,\a)}) S_\a(h_{(3,\a)}) \v_\a(h_{(4,\a)}) 1_e \\
&=& \v_\a(h_{(1,\a)}) \v_\a(h_{(2,\a)}) 1_e \\
&=& \v_\a(h)1_e.
\end{eqnarray*}

So $m_{\a,\bar{\a}}(id_{H_{\a}} \otimes T_\a)\Delta_\a = \v_\a 1_e$. Similarly prove $m_{\bar{\a},\a}(T_\a \otimes id_{H_{\a}})\Delta_\a = \v_\a 1_e$. 
Hence, $H$ with the new multiplication $m=\{m_{\a,\b}\}_{\a\in\pi}$ and antipode $T=\{T_{\a,\b}\}_{\a\in\pi})$  is a cocommutative Hopf $\pi$-algebra

(2) For all $h\in H_\a, ~\a\in\pi$, we have
\begin{eqnarray*}
& & B_\a(h_{(1,\a)})B_{\bar{\a}}(T_\a(h_{(2,\a)})) \\
&=& B_\a(h_{(4,\a)}) B_{\bar{\a}}(S_{\bar{\a}}(B_\a(h_{(2,\a)}))S_\a(h_{(3,\a)})B_\a(h_{(4,\a)}))  \\
&=& B_e(h_{(1,\a)}B_\a(h_{(2,\a)}) S_{\bar{\a}}(B_\a(h_{(3,\a)}))S_\a(h_{(4,\a)})B_\a(h_{(5,\a)}) S_{\bar{\a}}(B_\a(h_{(6,\a)}))) \\
&=& B_e(h_{(1,\a)} \v_\a(h_{(2,\a)}) S_\a(h_{(3,\a)}) \v_\a(h_{(4,\a)}) 1_e) \\
&=& \v_\a(h_{(1,\a)}) \v_\a(h_{(2,\a)}) B_e(1_e) \\
&=& \v_\a(h)1_e.
\end{eqnarray*}

From the above and $B_\a(h_{(1,\a)})S_{\bar{\a}}(B_\a(h_{(2,\a)}))=\v_\a(B_\a(h))=\v_\a(x)1_e$, we have $B_{\bar{\a}}T_\a$ and $S_{\bar{\a}}B_\a$ are the convolutional inverse for $B_\a$. Hence $B_{\bar{\a}}T_\a = S_{\bar{\a}}B_\a$, for all $\a\in\pi$.

(3) From Eq. (\ref{eq:RBO}) and (\ref{eq:RBHm}), we have 
\begin{eqnarray*}
 B_{\a\b}(g \circ_B h) &=& B_{\b\a}(g \circ_B h) \\
 &=& B_{\b\a}( g_{(1,\a)}B_\a(g_{(2,\a)})h S_{\bar{\a}}(B_\a(g_{(3,\a)})) ) \\
 &=& B_\a(g) B_\b(h),
\end{eqnarray*}
for all $g\in H_\a, h\in H_\b, ~\a,\b\in\pi$.

Then we have
\begin{eqnarray*}
& & B_{\b\a}\big(g_{(1,\a)}\circ_B B_\a(g_{(2,\a)})\circ_B h \circ_B T_{\bar{\a}}(B_\a(g_{(3,\a)}))\big) \\
&=& B_{\a}(g_{(1,\a)}) B_{\bar{\a}}(B_\a(g_{(2,\a)})) B_{\b}(h) B_{\a}(T_{\bar{\a}}(B_\a(g_{(3,\a)}))) \\
&=& B_{\a}(g_{(1,\a)}) B_{\bar{\a}}(B_\a(g_{(2,\a)})) B_{\b}(h) S_{\a}(B_{\bar{\a}}(B_\a(g_{(3,\a)})))  \\
&=& B_{\a}(g) \circ_B B_{\a}(h).
\end{eqnarray*}
for all $g\in H_\a, h\in H_\b, ~\a,\b\in\pi$. The operator $B$ is also a Rota-Baxter operator on the cocommutative Hopf $\pi$-algebra $H_B$.

(4) From (3), it is evident that the family of maps $B:H_B\rightarrow H$ is a homomorphism of Rota-Baxter Hopf $\pi$-algebras.
\end{proof}

\end{document}
